%!TeX root=../monografia
%("dica" para o editor de texto: este arquivo é parte de um documento maior)
% para saber mais: https://tex.stackexchange.com/q/78101
\chapter{Rascunho VEB}
\label{cap:rascunho-veb}

\section{Introdução}

\subsection{Problema do predecessor}

O problema do predecessor tem como objetivo manter um conjunto $S$ de $n$
elementos de um universo $U$ sujeito às seguintes operações:
\[
\text{Predecessor}(x) = \max \{y \in S \mid y < x \} 
\]
\[
\text{Sucessor}(x) = \min \{y \in S \mid y > x \}
\]
%(Mauricio): x \in U nas consultas também? Imagino que sim.

No nosso caso, estamos interessados no problema dinâmico, que também suporta as
operações de \textbf{Inserção}, para $x \in U$, e \textbf{Remoção}, para $x \in
S$.

Com uma Árvore Binária de Busca Balanceada (ABBB), conseguimos resolver o
problema com custo $\Theta (\lg n)$ para as operações e ocupando espaço $\Theta
(n)$.
%(Mauricio): Aqui você já está supondo Word RAM, visto que está considerando que os ponteiros de uma AABB têm tamanho O(1) cada.
No modelo clássico baseado em comparações, as únicas operações entre os valores
dos elementos são de comparação (como $\geq, =, \leq$) e nele uma ABBB é
assintoticamente ótima.

\subsection{Modelo Word RAM}

Vamos usar o modelo \textbf{word RAM} (Random Access Machine), definido por
Fredman e Willard em 1983 [conferir veracidade]. Nele, a memória é dividida em
\textit{palavras} com tamanho de $w$ bits. Assim, todos os elementos são
inteiros de $w$ bits e definimos um universo finito $U=\{0,1,\dots,u-1\}$, com
$u=2^w$.  

As operações aritméticas e lógicas entre elementos são executadas em tempo
constante:
\[+, -, *, /,\%,<<,>>,\&,|,\geq,=,\leq\]

Além disso, temos que $w\geq \lg n$, caso contrário, nem seria possível
representar o problema nesse modelo. [desenvolver]
%(Mauricio): Acho que $n$ deveria ser $u$.

Também vale notar que esse modelo já se aproxima da arquitetura real dos
computadores hoje em dia. Atualmente, um sistema comum tem $64$ bits, ou seja,
suportaria um universo de palavras de $0$ a $2^{64}$.

\section{Árvore de van Emde Boas}

\subsection{Introdução}

A estrutura foi inventada por um time liderado por Peter van Emde Boas em
$1976$. Originalmente, a Árvore de van Emde Boas (vEB) é proposta como uma
espécie de fila de prioridade melhorada, que usa o que eles chamam de ``árvores
estratificada''. Aqui, vou descrever em mais detalhes outra versão da estrutura,
com arrays [onde foi proposta pela primeira vez?], que consegue fazer as
operações descritas anteriormente em tempo $O(\lg \lg u)$ com espaço
$\Theta(u)$.

Uma ideia simples para lidar com o problema do predecessor seria construir um
\textit{bit vector}, isto é, um array com o mesmo tamanho do universo, $u$, e
para cada elemento, o valor $0$ no vetor indica que ele não está em $S$, e $1$
indica que está. Por exemplo, para um universo com $u=16$ e $S=\{2,3,5,7,11\}$
temos um \textit{bit vector} $A$ que representamos na
Figura~\ref{f:veb-bit-vector}.

\begin{figure}[h]
\centering
\begin{tikzpicture} 
    \drawBitArray[A]{0,0,1,1,0,1,0,1,0,0,0,1,0,0,0,0}{0,0}
\end{tikzpicture}
\caption{\textit{Bit vector} $A$ representando o conjunto $S={2,3,5,7,11}$
dentro de um universo com $u=16$. \label{f:veb-bit-vector}}
\end{figure}

Isso é muito eficiente para consultar se um elemento $x$ pertence a $S$, pois
basta checar o valor do array no índice $x$, em tempo constante. Contudo, para
consultas de predecessor o custo de pior caso é linear, $O(u)$. 

Para melhorar isso, uma ideia simples é construir uma árvore binária em cima do
vetor. Os elementos do vetor são as folhas e, de baixo para cima, definimos
sucessivamente os pais de cada nó, fazendo o \textbf{OR} entre seus filhos até
chegar na raiz. Cada nó terá valor 1 ou 0, indicando se pelo menos um dos seus
sucessores vale $1$. No final, obtemos uma estrutura que pode ser visualizada
com a Figura~\ref{f:veb-bst-bit-vector}

\enlargethispage{3\baselineskip}

\begin{figure}[!htbp]
\centering
    \begin{tikzpicture}
    % 1. Array Base A
    \drawBitArray[A]{0,0,1,1,0,1,0,1,0,0,0,1,0,0,0,0}{0,0}

    % 2. Camada 3 (folhas da árvore - combinações de 2 em 2)
    \foreach \i [evaluate=\i as \j using int(2*\i), evaluate=\i as \k using int(2*\i+1)] in {0,...,7} {
        \path ($(A\j.north)!.5!(A\k.north)$) ++(0, 0.9) node[veb node] (l3_\i) {};
        \draw (l3_\i) -- (A\j.north);
        \draw (l3_\i) -- (A\k.north);
    }
    % Valores da Camada 3
    \foreach \i/\val in {0/0, 1/1, 2/1, 3/1, 4/0, 5/1, 6/0, 7/0} {
        \node at (l3_\i) {\val};
    }

    % 3. Camada 2
    \foreach \i [evaluate=\i as \j using int(2*\i), evaluate=\i as \k using int(2*\i+1)] in {0,...,3} {
        \path ($(l3_\j.north)!.5!(l3_\k.north)$) ++(0, 0.9) node[veb node] (l2_\i) {};
        \draw (l2_\i) -- (l3_\j);
        \draw (l2_\i) -- (l3_\k);
    }
    \foreach \i/\val in {0/1, 1/1, 2/1, 3/0} { \node at (l2_\i) {\val}; }

    % 4. Camada 1
    \foreach \i [evaluate=\i as \j using int(2*\i), evaluate=\i as \k using int(2*\i+1)] in {0,1} {
        \path ($(l2_\j.north)!.5!(l2_\k.north)$) ++(0, 0.9) node[veb node] (l1_\i) {};
        \draw (l1_\i) -- (l2_\j);
        \draw (l1_\i) -- (l2_\k);
    }
    \foreach \i/\val in {0/1, 1/1} { \node at (l1_\i) {\val}; }

    % 5. Raiz
    \path ($(l1_0.north)!.5!(l1_1.north)$) ++(0, 0.9) node[veb node] (root) {1};
    \draw (root) -- (l1_0);
    \draw (root) -- (l1_1);   
\end{tikzpicture}
\caption{A estrutura representa uma árvore binária construída em cima do
\textit{bit vector} $A$. Cada nó tem valor 1 ou 0, indicando se pelo menos um de
seus sucessores tem valor 1.\label{f:veb-bst-bit-vector}}
\end{figure}

Agora, para encontrar o predecessor de $x$ fazemos os seguintes passos:
Começamos na folha dada pelo índice $x$ do array, então subimos os nós em
direção a raiz. Paramos se chegarmos em um nó vindo do seu filho direito e esse
nó tiver um $1$ como seu filho esquerdo. Por fim, descemos pela subárvore desse
filho buscando seu valor máximo, ou seja, sempre indo pelo filho mais a direita
com valor $1$.

Tome, por exemplo, a busca pelo predecessor de $7$ na estrutura de antes. A
Figura~\ref{f:veb-bst-bit-vector-search} mostra o processo dessa busca, obtendo
que $\Predecessor(7)=5$.

\begin{figure}[h]
\centering
    \begin{tikzpicture}
    % 1. Array Base A
    \drawBitArray[A]{0,0,1,1,0,1,0,1,0,0,0,1,0,0,0,0}{0,0}

    % 2. Camada 3 (folhas da árvore - combinações de 2 em 2)
    \foreach \i [evaluate=\i as \j using int(2*\i), evaluate=\i as \k using int(2*\i+1)] in {0,...,7} {
        \path ($(A\j.north)!.5!(A\k.north)$) ++(0, 0.9) node[veb node] (l3_\i) {};
        \draw (l3_\i) -- (A\j.north);
        \draw (l3_\i) -- (A\k.north);
    }
    % Valores da Camada 3
    \foreach \i/\val in {0/0, 1/1, 2/1, 3/1, 4/0, 5/1, 6/0, 7/0} {
        \node at (l3_\i) {\val};
    }

    % 3. Camada 2
    \foreach \i [evaluate=\i as \j using int(2*\i), evaluate=\i as \k using int(2*\i+1)] in {0,...,3} {
        \path ($(l3_\j.north)!.5!(l3_\k.north)$) ++(0, 0.9) node[veb node] (l2_\i) {};
        \draw (l2_\i) -- (l3_\j);
        \draw (l2_\i) -- (l3_\k);
    }
    \foreach \i/\val in {0/1, 1/1, 2/1, 3/0} { \node at (l2_\i) {\val}; }

    % 4. Camada 1
    \foreach \i [evaluate=\i as \j using int(2*\i), evaluate=\i as \k using int(2*\i+1)] in {0,1} {
        \path ($(l2_\j.north)!.5!(l2_\k.north)$) ++(0, 0.9) node[veb node] (l1_\i) {};
        \draw (l1_\i) -- (l2_\j);
        \draw (l1_\i) -- (l2_\k);
    }
    \foreach \i/\val in {0/1, 1/1} { \node at (l1_\i) {\val}; }

    % 5. Raiz
    \path ($(l1_0.north)!.5!(l1_1.north)$) ++(0, 0.9) node[veb node] (root) {1};
    \draw (root) -- (l1_0);
    \draw (root) -- (l1_1);   

    % --- DESENHO DOS CAMINHOS / SETAS DA BUSCA ---
    % Seta paralela na descida (exemplo: buscando o 7)
    \draw[path arrow] ($(l2_1.south west)$) -- ($(l3_2.north)$);
    \draw[path arrow] ($(l3_2.south)$) -- ($(A5.north)+(-0.1,0.0)$);

    % Seta na subida (exemplo: do nó 14 até a raiz)
    \draw[path arrow] ($(A7.north)+(0.1,0.0)$) -- ($(l3_3.south)+(0.2,0.1)$);
    \draw[path arrow] ($(l3_3.north west)+(0.2,0.1)$) -- ($(l2_1.south)+(0.2,0.1)$);
    
\end{tikzpicture}
\caption{Processo de busca pelo predecessor de $7$ na estrutura da árvore
binária construída em cima de $A$. Encontramos que o predecessor de $7$ é o
elemento $5$.\label{f:veb-bst-bit-vector-search}}
\end{figure}

Como a altura da árvore é $\lg u$ e essa operação vai fazer no máximo uma
varredura até a raíz e depois até uma folha, a complexidade desse algoritmo é
$O(\lg u)$. Ela acaba sendo muito similar a uma ABBB comum para a maioria das
operações.

Agora, vamos expandir isso para integrar ao conceito de clusters. Podemos dizer,
que cada um dos nós da árvore define um cluster como sua subárvore, e seu valor
nos diz se há elementos ou não naquele cluster. Os nós logo acima das folhas tem
um cluster de tamanho $2$, acima deles, o pai define um cluster de tamanho $4$. 
%(Mauricio): A frase acima não me parece gramaticalmente correta.

Note que, isso é definido pelo grau da árvore, o número de filhos que cada nó
tem. Podemos mudar a construção da árvore para que ela tenha um grau fixo
diferente, como $\sqrt u$. No nosso exemplo, $\sqrt u = \sqrt{16} = 4$. Com
isso, obtemos uma árvore como a da Figura~\ref{f:veb-clusters}.

\enlargethispage{1\baselineskip}

\begin{figure}[h]
\centering
\begin{tikzpicture}

    % 1. Desenha o vetor A na base usando a macro genérica
    \drawBitArray[A]{0,0,1,1,0,1,0,1,0,0,0,1,0,0,0,0}{0,0}

    % 2. Camada intermediária (4 nós com valores 1, 1, 0, 1)
    \def\midvals{1,1,1,0}
    
    \foreach \val [count=\k from 0] in \midvals {
        \pgfmathtruncatemacro{\start}{\k*4}
        \pgfmathtruncatemacro{\finish}{\start+3}
        
        % Posiciona o nó centralizado sobre cada bloco de 4 elementos
        \path ($(A\start.north)!.5!(A\finish.north)$) ++(0, 1.3) node[veb node] (m\k) {\val};
        
        % Conecta às 4 células correspondentes do vetor A
        \foreach \i in {\start,...,\finish} {
            \draw (m\k) -- (A\i.north);
        }
    }

    % 3. Nó Raiz no topo
    \path ($(m0.north)!.5!(m3.north)$) ++(0, 1.3) node[veb node] (root) {1};

    % Conecta a raiz aos 4 nós intermediários
    \foreach \k in {0,1,2,3} {
        \draw (root) -- (m\k);
    }
\end{tikzpicture}
\caption{Árvore de grau $\sqrt u = \sqrt{16} = 4$ construída em cima do
\textit{bit vector} $A$. Cada nó interno define um cluster que diz se cada um
dos 4 elementos que ele reúne tem pelo menos um bit 1.\label{f:veb-clusters}}
\end{figure}

\textit{[explicar o motivo de escolher $\sqrt u$ ]} 
 
Na prática, como o grau da árvore pode ser arbitrariamente grande, é usual que
os filhos de cada nó sejam dispostos em um vetor. No caso da nossa estrutura, é
um pouco diferente, todas as folhas já fazem parte de um mesmo
\textit{bit-vector} $A$. Mas é interessante para nós colocar os filhos da raiz
em um mesmo vetor. Uma estrutura construída assim é visualizada com a
Figura~\ref{f:veb-clusters-vector}.

\begin{figure}[h]
\centering
\begin{tikzpicture}

    % 1. Array Base A
    \drawBitArray[A]{0,0,1,1,0,1,0,1,0,0,0,1,0,0,0,0}{0,0}

    % 2. Linhas verticais grossas de divisão de Clusters (a cada 4 elementos)
    \foreach \x in {3, 7, 11} {
        \draw[cluster split] ($(A\x.north east)+(0, 0.1)$) -- ($(A\x.south east)-(0, 0.1)$);
    }

    % 3. Braces (chaves) sobre os clusters
    \draw[cluster brace] (A0.north west) -- (A3.north east) 
        node[midway, above=4pt] (b0) {};
    \draw[cluster brace] (A4.north west) -- (A7.north east) 
        node[midway, above=4pt] (b1) {};
    \draw[cluster brace] (A8.north west) -- (A11.north east) 
        node[midway, above=4pt] (b2) {};
    \draw[cluster brace] (A12.north west) -- (A15.north east) 
        node[midway, above=4pt] (b3) {};

    % 4. Array de Summary
    % Posicionado acima do ponto médio do vetor A
    \path ($(A0.north)!.5!(A15.north)$) ++(-0.825, 2.0) coordinate (sum_pos);
    
    \foreach \val [count=\i from 0] in {1, 1, 1, 0} {
        \node[array cell] (sum\i) at ($(sum_pos) + (\i*0.55, 0)$) {\val};
        %\node[array index, above=of sum\i] {\i};
    }

    % 5. Conexões entre o Summary e os Braces
    \draw (sum0.south) -- (b0.north);
    \draw (sum1.south) -- (b1.north);
    \draw (sum2.south) -- (b2.north);
    \draw (sum3.south) -- (b3.north);

    % 6. Raiz do summary
    %\path ($(sum0.north)!.5!(sum3.north)$) ++(-0.825, 2.0) coordinate (raiz_pos);


    \path ($(sum0.north)!.5!(sum3.north)$) ++(0, 1.5) node[veb node] (root) {1};

    \draw[cluster brace] (sum0.north west) -- (sum3.north east) node[midway, above=6pt] (brace) {};

    % Conecta a raiz aos 4 nós intermediários
    \draw (root) -- (brace);

\end{tikzpicture}
\caption{Árvore de grau 4 construída em cima do \textit{bit-vector} $A$ com seus
nós internos reunidos em um vetor.\label{f:veb-clusters-vector}}
\end{figure}

\textit{[meio feia essa figura, não sei se gostei]}

A partir da raiz, conseguimos acessarmente direta o vetor de clusters. Com o
índice do cluster de interesse, podemos obter a informação se ele está vazio ou
não instantaneamente. 

Vamos nos aprofundar nessa divisão mais tarde, antes vamos ver como isso nos
permite definir um sistema de coordenadas hierárquico. Nele, um elemento $x \in
U$ pode ser escrito da forma $x=c\sqrt u + i$, ou seja, estamos dividindo $x$
por $\sqrt u$, e temos o $c$ como quociente e $i$ como resto. Podemos obter os
valores de $c$ e $i$ com algumas contas:
\[
c = \left \lfloor \frac{x}{\sqrt u} \right \rfloor 
\]
\[
i = x \mod \sqrt u
\]

É interessante pensar também na perspectiva binária, pois as operações com bits
em um computador real são feitas muito mais rápidas que as de divisão. Por conta
disso, é com essas operações que a implementação da vEB será feita. 

Temos que $x$ é uma palavra de $w$ bits e $u= 2^w$, assim $\sqrt u = \sqrt{2^w}
= 2^{\frac{w}{2}}$. Ou seja, dividir $x$ por $\sqrt u$ na realidade é fazer um
bit-shift de $\frac{w}{2}$ bits. Podemos visualizar isso com um exemplo simples,
com $w=4$ e $x=11$. A representação dos bits de $x$ é dada pela
Figura~\ref{f:veb-hierarch-coords}.

\begin{figure}[h]
\centering
\begin{tikzpicture}

    % 1. Bits de x
    \def\bitvals{1,0,1,1}
    \foreach \val [count=\i from 0] in \bitvals {
        \node[array cell] (x\i) at ($(0,0) + (\i*0.55, 0)$) {\val};
    }
    \node[left=0.2cm of x0] {$x$};
    
    % 2. Braces (chaves) definindo c e i
    \draw[cluster brace] (x0.north west) -- (x1.north east) 
        node[midway, above=4pt] (b0) {};
    \draw[cluster brace] (x2.north west) -- (x3.north east) 
        node[midway, above=4pt] (b1) {};

    % Rotulos
    \node[above right=0pt and -10pt of b0, font=\small] {$c$};
    \node[above right=0pt and -10pt of b1, font=\small] {$i$};

\end{tikzpicture}
\caption{O elemento $x=11$ na sua representação como \textit{bit string},
$1011$. Seus 2 bits mais à esquerda são $c$, enquanto $i$ é dado pelos mais à
direita.\label{f:veb-hierarch-coords}}
\end{figure}

Em termos de operações fazemos:
\begin{align*}
c & = x >> w/2 \\
i & = x \And ((1 << w/2)-1) \\
x & = (c << w/2) | i
\end{align*}

Adotamos a partir daqui a notação $x=\langle c,i \rangle$ para representar a
chamada de funções que realizam essas operações para obter algum dos elementos.
Todas são feitas em tempo constante.                        

Vamos voltar para a nossa estrutura que constrói uma árvore em cima do
\textit{bit vector} de grau fixo $\sqrt u$. Essa divisão nos dá $\sqrt u$
clusters de tamanho $\sqrt u$. Seguindo com nosso exemplo, temos $\sqrt u = w =
4$ e $S=\{2,3,5,7,11\}$. A Figura~\ref{f:veb-clusters-vector-x} para melhor
visualização oculta o nó raiz.

\begin{figure}[h]
\centering
\begin{tikzpicture}

    % 1. Desenha o vetor A na base
    \def\bitvals{0,0,1,1,0,1,0,1,0,0,0,1,0,0,0,0}
    \foreach \val [count=\i from 0] in \bitvals {
        \ifnum\i=11
            \node[array cell, fill=Lavender] (A\i) at ($(0,0) + (\i*0.55, 0)$) {\val};
        \else
            \node[array cell] (A\i) at ($(0,0) + (\i*0.55, 0)$) {\val};
        \fi
        \node[array index, below=of A\i] {\i};
    }
    \node[left=0.2cm of A0] {$A$};

    % 2. Linhas verticais grossas de divisão de Clusters (a cada 4 elementos)
    \foreach \x in {3, 7, 11} {
        \draw[cluster split] ($(A\x.north east)+(0, 0.1)$) -- ($(A\x.south east)-(0, 0.1)$);
    }

    % 3. Braces (chaves) sobre os clusters
    \draw[cluster brace] (A0.north west) -- (A3.north east) 
        node[midway, above=4pt] (b0) {};
    \draw[cluster brace] (A4.north west) -- (A7.north east) 
        node[midway, above=4pt] (b1) {};
    \draw[cluster brace] (A8.north west) -- (A11.north east) 
        node[midway, above=4pt] (b2) {};
    \draw[cluster brace] (A12.north west) -- (A15.north east) 
        node[midway, above=4pt] (b3) {};

    % 4. Array de Summary
    % Posicionado acima do ponto médio do vetor A
    \path ($(A0.north)!.5!(A15.north)$) ++(-0.825, 2.0) coordinate (sum_pos);
    
    \foreach \val [count=\i from 0] in {1, 1, 1, 0} {
        \ifnum\i=2
            \node[array cell,  fill=Orchid] (sum\i) at ($(sum_pos) + (\i*0.55, 0)$) {\val};
        \else
            \node[array cell] (sum\i) at ($(sum_pos) + (\i*0.55, 0)$) {\val};
        \fi
        \node[array index, above=of sum\i] {\i};
    }

    % 5. Conexões entre o Summary e os Braces
    \draw (sum0.south) -- (b0.north);
    \draw (sum1.south) -- (b1.north);
    \draw (sum2.south) -- (b2.north);
    \draw (sum3.south) -- (b3.north);

\end{tikzpicture}
\caption{Árvore de grau 4 construída em cima do \textit{bit-vector} $A$ com seus
nós internos reunidos em um vetor, o nó raiz é ocultado. O elemento $x=11$ é
decomposto nas coordenadas $c=2$ e $i=3$, pois $x=c\times \sqrt{16}+i$. O
cluster de índice $c=2$ está pintado de roxo. O elemento de índice $i=3$ dentro
do cluster, índice $x=11$ em $A$, é destacado em
rosa.\label{f:veb-clusters-vector-x}}
\end{figure}

Destacamos que o sistema hierárquico de coordenadas consegue localizar o
elemento $x=11$ de forma simples. Como $x=c\times \sqrt{16} + i$, obtemos que
$c=2$ e $i=3$. Ou seja, $x$ está no cluster $2$, destacado em roxo, e, dentro
desse cluster, possui índice $3$, esse elemento em rosa na figura. 

Conseguimos então, armazenar a representação binária de $x=1011$ em duas
palavras de tamanho $w/2$, pois $c=10$ e $i=11$.

% ==============================================================================
% FIGURA 2: Decomposição vEB (Summary, Clusters e Braces)
% ==============================================================================
%\begin{figure}[h]
%\centering
%\begin{tikzpicture}
%
%    % 1. Array Base A
%    \drawBitArray[A]{0,0,1,1,0,1,0,1,0,0,0,1,0,0,0,0}{0,0}
%
%    % 2. Linhas verticais grossas de divisão de Clusters (a cada 4 elementos)
%    \foreach \x in {3, 7, 11} {
%        \draw[cluster split] ($(A\x.north east)+(0, 0.1)$) -- ($(A\x.south east)-(0, 0.1)$);
%    }
%
%    % 3. Braces (chaves) sobre os clusters
%    \draw[cluster brace] (A0.north west) -- (A3.north east) 
%        node[midway, above=4pt] (b0) {};
%    \draw[cluster brace] (A4.north west) -- (A7.north east) 
%        node[midway, above=4pt] (b1) {};
%    \draw[cluster brace] (A8.north west) -- (A11.north east) 
%        node[midway, above=4pt] (b2) {};
%    \draw[cluster brace] (A12.north west) -- (A15.north east) 
%        node[midway, above=4pt] (b3) {};
%
%    % Rotulo \sqrt{u} bits
%    \node[above right=2pt and -10pt of b3, font=\small] {$\sqrt{u}\text{ bits}$};
%
%    % 4. Array de Summary
%    % Posicionado acima do ponto médio do vetor A
%    \path ($(A0.north)!.5!(A15.north)$) ++(-0.825, 2.0) coordinate (sum_pos);
%    
%    \foreach \val [count=\i from 0] in {1, 1, 1, 0} {
%        \node[array cell] (sum\i) at ($(sum_pos) + (\i*0.55, 0)$) {\val};
%        \node[array index, above=of sum\i] {\i};
%    }
%    \node[left=0.2cm of sum0, font=\itshape] {summary};
%    \node[right=0.2cm of sum3, font=\small] {$\sqrt{u}\text{ bits}$};
%
%    % 5. Conexões entre o Summary e os Braces
%    \draw (sum0.south) -- (b0.north);
%    \draw (sum1.south) -- (b1.north);
%    \draw (sum2.south) -- (b2.north);
%    \draw (sum3.south) -- (b3.north);
%
%\end{tikzpicture}
%\end{figure}
%

Podemos construir um algoritmo intuitivo para uma consulta
\textit{Predecessor}$(x)$ com o que temos até agora:

\begin{itemize}
    \item Procuramos dentro do cluster de $x$ pelo primeiro $1$ à sua esquerda;
    \item Se não achamos, procuramos pelo primeiro $1$ à esquerda de $c$, ou
    seja o primeiro cluster não vazio à esquerda do cluster de $x$; 
    \item Encontrado esse cluster, queremos o maior elemento dentro dele. O $1$
    mais à direita.
\end{itemize}

Essa rotina simples é mais eficiente que o \textit{bit vector}, mas cada passo
dela custa $O(\sqrt{u})$, então ela ainda é pior que aquela árvore binária que
construímos antes. 

Para melhorar isso, podemos expandir essa ideia recursivamente. Nossa estrutura
vEB terá $\sqrt{u}$ clusters e um sumário, e eles serão também estruturas vEB
contendo clusters e sumários.

No momento, uma árvore vEB denominada $V$ de tamanho $u$ com tamanho de
\textit{palavra} $w$ contém:

\begin{itemize}
\item $V.clusters:$ Array de $\sqrt u$ Árvores vEBs de tamanho $\sqrt u$ com
tamanho de \textit{palavra} $\frac{w}{2}$. 
\item $V.summary:$ Árvore vEB de tamanho $\sqrt{u}$ com tamanho de palavra
$\frac{w}{2}$
\end{itemize}

O Programa~\ref{p:veb-first-predecessor} aplica um algoritmo para encontrar o
predecessor recursivamente, seguindo a ideia da rotina anterior.

\begin{programruledcaption}{Função predecessor recursiva\label{p:veb-first-predecessor}}
  \begin{lstlisting}[
    language=pseudocode,
    style=pseudocode,
    style=wider,
    functions={Predecessor},
    specialidentifiers={},
  ]
    Function Predecessor$(V, x=\langle c,i \rangle)$
        $i$ := Predecessor$(V.cluster[c], i)$
        if $i$ == $-\infty$
            $c$ := Predecessor$(V.summary, c)$
            $i$ := Predecessor$(V.cluster[c], +\infty)$
        return $\langle c,i \rangle$
  \end{lstlisting}
\end{programruledcaption}

É interessante também pensar no algoritmo de inserção nessa estrutura, dado no
Programa~\ref{p:veb-first-insert}

\begin{programruledcaption}{Função recursiva de inserção \label{p:veb-first-insert}}
  \begin{lstlisting}[
    language=pseudocode,
    style=pseudocode,
    style=wider,
    functions={Inserir},
    specialidentifiers={},
  ]
    Function Inserir$(V,x=\langle c,i\rangle)$
        Inserir$(V.summary, c)$
        Inserir$(V.cluster[c], i)$
  \end{lstlisting}
\end{programruledcaption}

Lembrando que a operação $x=\langle c, i \rangle$ é feita em tempo constante, o
tempo de execução dessa rotina é dado pela recorrência:
\[
T(u) = 2T(\sqrt{u}) + O(1)
\]

Ou então,
\[
T'(\lg u) = 2T'\left (\frac{\lg u}{2}\right ) + O(1)
\]

Resolvendo, obtemos um custo de $O(\lg u)$.

Podemos também resolver a recorrência do \textit{Predecessor}, que tem
complexidade $O(\lg u^{\lg 3})$. \textit{[Provar essa recorrência?]} Ainda está
longe do que gostaríamos, isso porque no pior caso fazemos três chamadas
recursivas de \textit{Predecessor}.

A chamada na linha 4 pode ser traduzida como buscar o maior elemento dentro de
um cluster. Se nós armazenamos isso de antemão, com um elemento \textit{max},
reduzimos para apenas 2 chamadas recursivas no
Programa~\ref{p:veb-predecessor-max}

\begin{programruledcaption}{Função predecessor recursiva com \emph{max} \label{p:veb-predecessor-max}}
  \begin{lstlisting}[
    language=pseudocode,
    style=pseudocode,
    style=wider,
    functions={Predecessor},
    specialidentifiers={},
  ]
    Function Predecessor$(V, x=\langle c,i \rangle)$
        $i$ := Predecessor$(V.cluster[c], i)$
        if $i$ == $-\infty$
            $c$ := Predecessor$(V.summary, c)$
            $i$ := $V.cluster[c].max$
        return $\langle c,i \rangle$
  \end{lstlisting}
\end{programruledcaption}

A primeira chamada busca no cluster $c$ pelo predecessor de $i$, se ela falha
então buscamos pelo predecessor desse cluster na estrutura de sumário.
Gostaríamos poder já saber antes se essa busca dentro do cluster $c$ terá
sucesso, pois caso contrário podemos pular essa chamada e ir direto para a
outra. A maneira mais intuitiva de fazer isso é armazenando também o elemento
mínimo de cada estrutura. Se $i$ for maior que $V.cluster[c].min$, então a busca
dentro do cluster deve ser realizada pois encontraremos o predecessor lá. Caso
contrário, podemos partir para a busca de outro cluster, sem precisar varrer
$c$. O algoritmo de predecessor então fica da forma apresentada no
Programa~\ref{p:veb-predecessor-min-max}

\begin{programruledcaption}{Função predecessor recursiva com \emph{min} e \emph{max}\label{p:veb-predecessor-min-max}}
  \begin{lstlisting}[
    language=pseudocode,
    style=pseudocode,
    style=wider,
    functions={Predecessor},
    specialidentifiers={},
  ]
    Function Predecessor$(V, x=\langle c,i \rangle)$
        if $i$ > $V.cluster[c].min$
            $i$ := Predecessor$(V.cluster[c], i)$
        else
            $c$ := Predecessor$(V.summary, c)$
            $i$ := $V.cluster[c].max$
        return $\langle c,i \rangle$
  \end{lstlisting}
\end{programruledcaption}

Com isso, temos apenas uma chamada recursiva e o custo da recorrência é dado
por:
\[
T(u) = T(\sqrt u) + O(1)
\]

Ou então, aplicando log temos $T'(\lg u) = T'(\lg \sqrt u) + O(1)$, que em
termos de $w$ é:
\[
T'(w) = T'\left (\frac{w}{2}\right ) + O(1)
\]

Isso, é claro, é resolvido para $O(\lg w)$, ou seja, $O(\lg \lg u)$, como
queríamos.

Com essa adição, o algoritmo de inserção ainda permanece com a mesma
complexidade, apenas mantendo registro do mínimo e máximo. O
Programa~\ref{p:veb-insert-min-max} aplica esse algoritmo.

\begin{programruledcaption}{Função recursiva de inserção com \emph{min} e \emph{max} \label{p:veb-insert-min-max}}
  \begin{lstlisting}[
    language=pseudocode,
    style=pseudocode,
    style=wider,
    functions={Inserir},
    specialidentifiers={},
  ]
    Function Inserir$(V,x=\langle c,i\rangle)$
        if $x$ < $V.min$
            $V.min$ := $x$
        if $x$ > $V.max$
            $V.max$ := $x$
        Inserir$(V.summary, c)$
        Inserir$(V.cluster[c], i)$
  \end{lstlisting}
\end{programruledcaption}

Para também obtermos a complexidade de $O(\lg \lg u)$ na inserção, queremos
poder fazer apenas uma das chamadas recursivas também. Note que a inserção na
estrutura de sumário que fazemos na linha 5 só é necessária quando o primeiro
elemento daquele cluster é inserido. A ideia é que nesse caso podemos fazer
apenas essa inserção, e não inserimos o elemento recursivamente no cluster,
apenas o salvamos a parte. O elemento escolhido para isso será o \textit{min},
que também servirá como indicativo se um cluster está vazio ou não. As funções
\textsc{empty}$(V)$ e \textsc{NotEmpty}$(V)$ vão retornar se $V.min$ tem o valor
especial $none$ ou não. 

O algoritmo fica da forma do Programa~\ref{p:veb-final-insert}.

\begin{programruledcaption}{Função recursiva de inserção com \emph{lazy min} \label{p:veb-final-insert}}
  \begin{lstlisting}[
    language=pseudocode,
    style=pseudocode,
    style=wider,
    functions={Inserir},
    specialidentifiers={},
  ]
    Function Inserir$(V,x=\langle c,i\rangle)$
        if $V.min$ = none
            $V.min$ := $V.max$ := $x$
            return
        if $x$ <= $V.min$
            \func{Swap}$(x,V.min)$
        if $x$ > $V.max$
            $V.max$ := $x$
        if \func{Empty}$(V.cluster[c])$
            Inserir$(V.summary, c)$
        Inserir$(V.cluster[c], i)$
  \end{lstlisting}
\end{programruledcaption}

Se a estrutura está vazia no caso base da primeira linha, apenas salvamos $x$
como o \textit{min} (e o \textit{max}) e retornamos sem mais chamadas. Então,
quando é feita a chamada na linha 8 de custo $O(\lg \lg u)$, sabemos que a
chamada da linha 10 terá custo constante pois cairá nesse caso. Assim, a
complexidade obtida é $O(\lg \lg u)$.

Vale notar que no caso da linha 4, o elemento que queremos inserir é menor que o
\textit{min} e, portanto, precisamos inserir o antigo \textit{min} na estrutura,
então é feito um \textit{swap} entre os valores.  

Com essas mudanças, precisamos fazer pequenas alterações no nosso código do
\textit{predecessor} para lidar com o caso especial do \textit{min} e das
estruturas vazias, resultando no Programa~\ref{p:veb-final-predecessor}

\begin{programruledcaption}{Função predecessor recursiva com \emph{lazy min} \label{p:veb-final-predecessor}}
  \begin{lstlisting}[
    language=pseudocode,
    style=pseudocode,
    style=wider,
    functions={Predecessor},
    specialidentifiers={},
  ]
    Function Predecessor$(V, x=\langle c,i \rangle)$
        if $x$ <= $V.min$
            return null
        if \func{NotEmpty}$(V.cluster[c])$ and $i$ > $V.cluster[c].min$
            return $\langle c,$ Predecessor$(V.cluster[c], i) \rangle$
        else
            $c'$ := Predecessor$(V.summary, c)$
            return $\langle c',V.cluster[c'].max \rangle$
  \end{lstlisting}
\end{programruledcaption}

\subsection{Estruturas e códigos finais largados abaixo}

Por fim, a estrutura final de uma árvore vEB denominada $V$ de tamanho $u$ com
tamanho de \textit{palavra} $w$ contém:
\begin{itemize}
\item $V.min:$ Palavra armazenada pela árvore em si.
\item $V.clusters:$ Array de $\sqrt u$ Árvores vEBs de tamanho $\sqrt u$ com
tamanho de \textit{palavra} $\frac{w}{2}$. 
\item $V.summary:$ Árvore vEB de tamanho $\sqrt{u}$ com tamanho de palavra
$\frac{w}{2}$
\item $V.max:$ Cópia do elemento máximo da árvore
\end{itemize}

Pseudocódigo das operações:

\begin{programruledcaption}{Função predecessor recursiva com \emph{lazy min} \label{p:veb-final-predecessor-1}}
  \begin{lstlisting}[
    language=pseudocode,
    style=pseudocode,
    style=wider,
    functions={Predecessor},
    specialidentifiers={},
  ]
    Function Predecessor$(V, x=\langle c,i \rangle)$
        if $x$ <= $V.min$
            return null
        if \func{NotEmpty}$(V.cluster[c])$ and $i$ > $V.cluster[c].min$
            return $\langle c,$ Predecessor$(V.cluster[c], i) \rangle$
        else
            $c'$ := Predecessor$(V.summary, c)$
            return $\langle c',V.cluster[c'].max \rangle$
  \end{lstlisting}
\end{programruledcaption}

\begin{programruledcaption}{Função sucessor recursiva com \emph{lazy min} \label{p:veb-final-successor-1}}
  \begin{lstlisting}[
    language=pseudocode,
    style=pseudocode,
    style=wider,
    functions={Sucessor},
    specialidentifiers={},
  ]
    Function Sucessor$(V, x=\langle c,i \rangle)$
        if $x$ <= $V.min$
            return $V.min$
        if \func{NotEmpty}$(V.cluster[c])$ and $i$ < $V.cluster[c].max$
            return $\langle c,$ Sucessor$(V.cluster[c], i) \rangle$
        else
            $c'$ := Predecessor$(V.summary, c)$
            return $\langle c',V.cluster[c'].min \rangle$
  \end{lstlisting}
\end{programruledcaption}

\begin{programruledcaption}{Função recursiva de inserção com \emph{lazy min} \label{p:veb-final-insert-1}}
  \begin{lstlisting}[
    language=pseudocode,
    style=pseudocode,
    style=wider,
    functions={Inserir},
    specialidentifiers={},
  ]
    Function Inserir$(V,x=\langle c,i\rangle)$
        if $V.min$ = none
            $V.min$ := $V.max$ := $x$
            return
        if $x$ <= $V.min$
            \func{Swap}$(x,V.min)$
        if $x$ > $V.max$
            $V.max$ := $x$
        if \func{Empty}$(V.cluster[c])$
            Inserir$(V.summary, c)$
        Inserir$(V.cluster[c], i)$
  \end{lstlisting}
\end{programruledcaption}

\begin{programruledcaption}{Função recursiva de remoção com \emph{lazy min} \label{p:veb-final-remove}}
  \begin{lstlisting}[
    language=pseudocode,
    style=pseudocode,
    style=wider,
    functions={Remover},
    specialidentifiers={},
  ]
    Function Remover$(V,x=\langle c,i\rangle)$
        if $x$ = $V.min$
            $c$ := $V.summary.min$
            if $c$ = none 
                $V.min$ := none
                return
        $x$ := $V.min$ := $\langle c, i = V.cluster[c].min \rangle$
        Remover$(V.cluster[c],i)$
        if $V.cluster[c].min$ = none
            Remover$(V.summary,c)$
        if $V.summary.min$ = none
            $V.max$ := $V.min$
        else 
            $c'$ := V.summary.max
            $V.max$ := $\langle c', V.cluster[c'].max \rangle$
  \end{lstlisting}
\end{programruledcaption}

Anotações feedback vEB:
\begin{itemize}
    \item Questão do espaço - esqueci de colocar o básico desse espaço $O(u)$ e falar da versão com espaço linear (hash tables e gancho para X/Y-fast tries)
    \item Pensar em exemplos melhores, mais didáticos - na busca por exemplo mostrar algum caso em que mais de um elemento do summary seja ``atravessado''  (como um summary = [1, 1, 0, 1] para uma busca de predecessor de um elemento no ultimo cluster (12,13,14,15))
    \item Mostrar as operações acontecendo por meio de figuras
\end{itemize}